

\chapter{IDEAS}
This chapter is here to leave place for unordered ideas

\section{Use a NN to evaluate gp-trees and train them to perform better?}

\section{Use GP to make neurons superfluous? Use NN to get these last improvements for the gp solution?}

\section{Function-based backpropagation}
Similar to backpropagation in neural networks, it is conceivable that a similar approach can be used for computational trees to evaluate the influence of nodes, whether positive or negative. This is done without deriving the functions.
For each node the influence is calculated as a percentage of each input.

Function: If:

condition: 0.5
Case (selected): 0.5
Case (not selected): 0

\smallskip
Example: for x = 3\\
if x < 0 then 0 else 2:\\
input 1: 0.5\\
input 2: 0\\
input 3: 0.5 \\
\smallskip
Function "+":\\
input 1: value 1 / (value 1 + value 2)\\
input 2: value 2 / (value 1 + value 2)\\
Example: 5 + 20\\
Share of input 1: 0.2\\
Share of input 2: 0.8\\

\smallskip
Function "min":\\
input 1: 0.25\\
input 2: 0.25\\
input that is "min": 0.5\\
Example: min(1, 2)\\
input 1: 0.75\\
input 2: 0.25\\
\smallskip
The right and wrong cases are considered separately. The result then has to be analyzed respectively.\\
If, for example, it turns out that case ``a'' of an if-function always chosen for right decisions, but never chosen in wrong decisions, its right value would be very high and its wrong value very low. This function can stay and case b has to be optimized.

Depending on the distribution of positive and negative cases, the same applies to the condition. If its value for positive cases is greater than its negative value, it very often successfully excludes the incorrect condition.
In summary, the input that has the greatest difference between negative cases and positive cases should be given priority.


\section{random... can this approach work on delicate envs?}
To efficiently use this approach for delicate environments, the following strategy was considered.
\begin{enumerate}
    \item Let a neuronal network copy the strategy of the heuristic.
    \item let the network improve itself in a way, that it can not generate too much harm. (How, though, should that work)
    \item Use genetic programming to improve the heuristic to the level of the neuronal network by genetic programming.
        
    \item problem-analysed changes\\
    When one capsuled part of the program is already found to do good improvement, genetic changes and improvements should specifically aim to optimize this part with genetic programming.
    \item modifiable nodes erweitern. wenn in allen kandidaten etwas gefunden wurde, einfach diesen kern erweitern. candidate mining.
    \item \url{https://en.wikipedia.org/wiki/Genetic_improvement_(computer_science)}
    \item machine learning leads to opportunities to interact more with humans like speech recognition, image recognition (race recognition), interpret gestures and many more. It is weird that while merging more and more towards human they lose their explainability
    \item any black box goes. just train it, we check it afterwards
    \item rate of finding fit or explainable candidates based on input factors
    \item kleineres neuronales netzwerk, dass als zusätzlichen input die Berechnungen von der gp bekommt?
    \item By using machine learning approaches, the decision to sacrifices explainability for performance improvement is made.
    \item Schaut man einen Baum an, kann man nicht von der Position eines Knotens direkt auf seine Relevanz schließen. Die Auswirkungen von kleinen Änderungen an Blattknoten ist dafür zu unvorhersehbar.
    \item If the basic structure of a program is broken up in an unexpected way, this is less a deterioration of the explainability than a discovery of a potential core problem of the existing logic.
    \item Neural networks take longer to evaluate than a program!
\end{enumerate}

\subsection{Complexest programs}

Due to the fact, that we  can also use a large program as origin, it is very likely that we can find large, complex programs that could never be found if we had started it from scratch.

\subsection{loops}
While- and for loops are not part of the genetic programming operators. It might be interesting for the developer to give the opportunity at some places to make specific loop operations optional. E. g. when a loop of specific size is expected.

\subsection{unknown observations in Neuronal Networks}
what happens, when the SARSA occurs something it has never seen at all before

\subsection{increasing complexity over time}


\subsection{complexity is amount of variables?}
This means, a super large program that only uses a single input once is very explainable
Also:
delete all parts of the tree which do not have inputs and then regard the size

\subsection{increased complexity over time}
Komplexität Grenze über Zeit erhöhen? Mehr SIcherheit bezüglich der einfachen Kandidaten

\subsection{data samples}
use regular samples, do not specifically choose seldom ones more often. sample size 2.000 is ok.

\subsection{unknown states}
Problem with using just data sets: some states may never be reached if we have other decisions beforehand.

\subsection{mutation types}
- point mutation
- filters (gauss 10\%, poisson?)
- branch mutations

\begin{itemize}
    \item no origin reproduction as reproduction. TED must be enough
    \item CMA-ES algorithm for the best operations:
    Statistik machen, protokollieren, welche veränderung am besten war… mit optimieren?
    \item Overfitting within the neuronal network is NOT a problem, as long as the data samples are not containing cases, that completely fail.
    \item in gp, usually a function is approached where the result is rounded to the actual action space. we give the opportunity to do this differently
    \item how important are output dimensions? what, if there are 1000?

\end{itemize}


\subsection{Discarded GP improvements}

Rating operations:\\
Some operations are more intuitive than others.

\begin{tabular}{l l}
     Addition & 0.5 \\
     Multiplication & 2\\
     Multiplication with values $\approx 1$ & 1\\
     If-then-else & 3\\
     
\end{tabular}

Note, that in order to use a weighting for the tree edit distance, the weighted tree edit distance must be used.

\section{Only gp-search unsolved examples}

\begin{itemize}
    \item cluster the data from the blackbox-samples that are not explainable with experts approach
    \item gp process should not affect the samples that are already correct 
\end{itemize}
